Item (4) is the capitalized annual expenditure of 900 thalers from which the variable part of the widow's pension is to be met. This sum is divided equally among all eligible widows if their number is 18 or more; if there are fewer, each receives 50 thalers. It follows that in the latter case (which also applies at this moment) the fund makes a saving that has not been included in the balance account and cannot be calculated in advance at all for a somewhat longer period. In any case, this is only a temporary advantage which, later on, once the effects of the association's present great expansion have developed, will occur rarely or perhaps never again.

\Needspace{8\baselineskip}
I have assumed the following average values for incidental expenditure
\begin{center}
\small\setlength{\tabcolsep}{2pt}
\begin{tabular}{r r l r l r l p{196pt}}
&28&Thl.&4&Ggr&—&Pf.&for legal costs, insofar as not reimbursed,\\
&61&,,&8&,,&9&,,&for building costs\\
&133&,,&9&,,&5&,,&for keeping the accounts and copying fees\\
&136&,,&6&,,&3&,,&for losses\\
Total&359&Thl.&4&Ggr.&5&Pf.&
\end{tabular}
\end{center}
\noindent\hspace*{\parindent}These estimates are based on the experience of the last 21 or 20 years. The calculation of 1845 could only be based on experience from 15 or 14 years, which had given the total of 323 Thl. 13 Ggr. The former figure, capitalized, yields 10262 Thl. 10 Ggr. This amount, with the addition of 2644 Thl. 4 Ggr. (as the mean value of the non-earning part of the assets over a 17-year average), forms item (5). For the latter figure, incidentally, the calculation of 1845, using an 11-year average, had assumed the nearly equal value of 2635 Thl. 18 Ggr. 4 Pf.

Item (6) arises from capitalizing the annual income from the lease rent of the university pharmacy (1000 thalers).
